% ======================= 第 六 章 =======================
\chapter{有穷氏遗迹与操作表}
\chaptersubtitle{ГЛАВА ШЕСТАЯ}{МАШИНА РОДА ЮЦИН}
\epigraphbox{“多行不义，必自毙。”}{《左传·隐公元年》}

\kaishou{本章只写一件操作：一个人照着表做了三步，没有做第四步。全书对她的动机不给一个字，档案只有一行「未解」。}

\noindent
前五章说的都是这座城的内部事务：制度、造物、规程、历法。这一章要走出城一次，去看地面上一处遗迹，以及一个人在那里做了什么。

这件事在传世文献里有名字，记作一次复仇。本书要说明的是，所谓复仇，在操作上是一次没有执行完的程序。

\section{地下构筑物的形制}

有穷氏的地望，传统说法在今山东平原一带，《左传》两处提到它，都与后羿有关。\footnote{襄公四年与昭公二十八年，皆见杨伯峻注本。两处对后羿的褒贬不同，注家已有讨论。}一九七〇年代以后，这一带陆续有商周遗存出土，与传说中的年代相去较远，两者能否对应，学界没有定论。

`ФВ-1972-06` 记的是另一样东西。它是一份地下构筑物测绘记录，出土地点在该区域的北部，深度 11.4 米，顶板为整块石板，接缝三处，均以铅锡封死。

顶板之下是一间长方形的室，长 6.8 米，宽 3.2 米，高 2.4 米。室内没有器物，四壁有刻线，刻线的走向与月面那套记号一致。

本章要用的是这间室的一处细节：室壁内侧，靠近出入口的一面，有一块金属板，尺寸 297 mm × 210 mm，厚 1.6 mm。板上刻着一张表。

\section{室壁操作表的结构}

表分六列，十四行。

左边五列是操作：范围、起点、方向、深度、终止。每一列下面有若干刻格，格内是记号。第六列在最右边，与左边五列之间隔着一道竖线。

第六列只有三个格，格内的记号与前五列不是同一套。

笔者请人做过复原，结论是：前五列是执行序列，第六列是退出序列。左五列决定做什么，最右一列决定怎么停。

这张表的排法本身就有问题。任何一份操作表，退出程序都会与主程序并列，或放在开头，因为人在紧急时要一眼找到它。这一份把它放在最远的一列，与主程序之间还隔着一道线。

\begin{figure}[htbp]
\centering
\begin{tikzpicture}[x=1mm,y=1mm]
  \useasboundingbox (-4,-16) rectangle (128,14);
  % 表框
  \draw[ink,line width=0.5pt] (0,-12) rectangle (116,8);
  % 竖线：五列 + 一列
  \foreach \x in {18,36,54,72,90}{\draw[ink,line width=0.3pt] (\x,-12) -- (\x,8);}
  \draw[ink,line width=0.8pt] (90,-12) -- (90,8);
  % 表头
  \foreach \x/\t in {9/范围, 27/起点, 45/方向, 63/深度, 81/终止, 103/退出}{
    \node[font=\scriptsize] at (\x,5.6) {\t};}
  % 记号：左五列有，右一列只有三格
  \foreach \y in {2.4,0.2,-2.0,-4.2,-6.4,-8.6}{
    \foreach \x/\s in {9/1, 27/0, 45/1, 63/1, 81/0}{
      \ifnum\s=1
        \fill[ink] (\x-2,\y) rectangle (\x+2,\y+0.9);
      \fi
    }
  }
  \foreach \y in {2.4,-2.0,-6.4}{
    \draw[ink,line width=0.3pt] (99,\y) rectangle (107,\y+0.9);}
  % 退出列虚线框
  \draw[ink,line width=0.6pt,dash pattern=on 2pt off 1.6pt] (92,-11) rectangle (114,3.8);
  \node[anchor=north,font=\scriptsize,align=center] at (103,-12.4) {退出序列};
  \node[anchor=north,font=\scriptsize,align=center] at (45,-13.4) {执行序列};
\end{tikzpicture}
\caption{`ФВ-1972-06` 室壁操作表的复原（局部）}
\figsource{据该档案的拓片与红外成像重绘。左侧实心块表示刻格已被磨平，记号无法辨读；只保留格位。}
\label{fig:table}
\end{figure}

这份表的用处，本来是给会读的人看的。问题在于，用它的人不识字。

\section{三次操作与缺失的第四步}

档案记载的操作只有三步。第一步，划定范围。范围以操作位置为原点，半径约四十步。这一步的记号在最左一列，操作者需要一个动作。

第二步，启动。这一步没有记号，只有一个凹槽，按下去即可。接连着的第三步，档案里没有。

严格地说，第二步启动之后，设备会自己走完前五列。操作者要做的是在执行完之后按第六列退出。她没有按。

`ФВ-1972-06` 附有一份现场记录，写的是事后三十年清理时的所见：室内四壁的刻线在某一高度处中断，中断处之上与之下材质不同，断面整齐。记录者的原话是「如被切开」。

这一段是全书最难写的地方。要说明它，只能回到第二章那个结论：她认得那些记号，因为幼年见过；她也能照着排下来的顺序做，因为旧历抄本给了她顺序。但她不识字，也不懂术数。一张表对她说来，是三个可以照做的动作，不是一个可以读懂的程序。

她做到了第三步该做的地方，只是不知道还有第四步。

\section{界线外移的五十年}

设备没有停在原处。按第一章所述那套装置的性质，被处理的范围不属于原来的整体，机体的修复程序没有可以作用的对象。这个范围不会自己合拢，也不会自己扩大；扩大的原因是另一边还在长。

边缘的组织仍在更新。每一次更新，新生的部分都以那一片为参照，于是被划进去。范围就这样一圈一圈往外走，起初只有一小块，五十年后走了样。

档案记她本人在范围内，一并被处理，当场死。这一条与「范围以操作者为原点」是一致的：她站在原点，没有退出。

至于她为什么站在那里，本书不给答案。档案只有一行，末尾注「未解」。笔者能确认的只有一件事：那一刻她并不需要站在原点，因为启动之后设备会自己走完。她站进去了。

\section{本章结论}

本章说了四件事。

遗迹是一间没有器物的室，四壁有与月面同源的刻线。室壁上有一张六列的操作表。

表的第六列是退出程序，位置最远，与主程序隔着一道线。这份排法不利于紧急操作。

操作只做了三步。前三步都能靠模仿完成，第四步需要读懂。

本人当场死，此后那台设备一直在走。

四件事合起来，与本书的总题目对得上：一个把「变化」删掉的文明，最后被一次没有执行完的程序处理掉；而执行它的人，是一个连字都不识的土著。

这句话容易被读成报应。笔者不这么看。它更像一次接口不匹配：造东西的人按自己的读法排了表，用东西的人按自己的读法用了它，两边都没有错，错在两套读法之间没有过渡。

本章的自我拆台：笔者不知道她为什么站在原点。不知道，就不写。全书关于动机的字数是零，这一节也一样。
